\documentclass[10pt,a4paper]{amsart}
\usepackage[margin=19mm]{geometry}
\usepackage{amsmath,amssymb,mathtools}
\usepackage[scr=rsfs,cal=euler]{mathalfa}
\usepackage{xfrac}
\usepackage[unicode,hidelinks,bookmarks=false]{hyperref}
\newcommand{\ie}{i.e.\ }
\newcommand{\cf}{cf.\ }
\newcommand{\resp}{respectively\ }
\newcommand{\Subsection}{\subsection}
\providecommand{\nobreakdash}{\nobreak\hbox{-}}
\setlength{\emergencystretch}{2em}
\multlinegap0pt
\usepackage{tikz}
\usepackage{tcolorbox}
\newtheorem{assumption}{Assumption}
\newtheorem{lemma}{Lemma}
\newcommand{\diag}{\mathrm{diag}}
\newcommand{\srank}{\mathrm{srank}}
\newcommand{\tr}{\mathrm{tr}}
\title{Rank-Efficient LoRA via Joint Tangent-Space Optimization under Isotropic Curvature}
\author{Zihan Zhu, Zhehang Du, Xuyang Chen, Tim Tsz-Kit Lau, Jiayuan Wu, X. Y. Han, Qi Long, Weijie Su}
\date{Source: arXiv:2609.12123v1 (CC BY 4.0)}
\begin{document}
\maketitle

\noindent\emph{Source and licence.} Rank-Efficient LoRA via Joint Tangent-Space Optimization under Isotropic Curvature. Version arXiv:2609.12123v1 (CC BY 4.0), distributed by arXiv under the Creative Commons Attribution 4.0 International licence, http://creativecommons.org/licenses/by/4.0/, as recorded in the arXiv licence metadata for this exact version and shown on the abstract page https://arxiv.org/abs/2609.12123v1. Full licence notice: \texttt{licenses/arxiv-2609.12123-CC-BY-4.0.txt}.\par\smallskip

\noindent\textit{Editorial scope.} Counted unit: the article's lemma lem:firststep-gd-rev (source0.tex lines 2235--2257). The article's complete original proof of it is delivered with it (source0.tex lines 2259--2421, one \texttt{proof} environment of 163 lines). No mathematics is written, repaired or re-derived: the statement and the proof are the article's own lines, in the article's own notation, and no retained line is substituted or re-flowed.  Retained uncounted as the source-local context the proof needs: lines 648--654; lines 732--741; lines 1936--1956.  Nothing was omitted from any retained span, so the package carries no omission record.  No retained line contains a citation, so the wrapper carries no bibliography and no bibliography entry.  No retained line contains a figure environment, an image-inclusion command or drawing code, so no image is delivered and the package declares no asset dependency.  Editorial formatting only, in addition: an AMS \texttt{amsart} preamble that restates the article's own declarations and macros used by the retained lines, namely the environments \texttt{assumption}, \texttt{lemma} on separate counters, together with the standard packages those lines need.  The retained environments print in this document as Assumption 1, Lemma 1 in order of appearance, whereas the article prints the same items with the numbering of its own full document.  The complete native source of the article as distributed in the arXiv e-print for this exact version is bundled unchanged as \texttt{sources/210080/source0.tex}.
\begin{equation}
\label{eq:init-firststep-rev}
(A_0)_{ij}\overset{\mathrm{i.i.d.}}{\sim}\mathcal N(0,1/m),
\qquad i\in[m],\; j\in[r],
\qquad
B_0=0 .
\end{equation}

\begin{assumption}[Rank-$k$ spiked gradient at initialization]
\label{assump:spike-init-rev}
There exist an integer $k\ge1$, a scalar $\zeta_0\ge0$, an orthonormal matrix $V_\star\in\mathbb R^{m\times k}$, and a diagonal matrix $\Theta_\star=\diag(\theta_1,\ldots,\theta_k)$ where $ \theta_1\ge\cdots\ge\theta_k>0$, such that
\begin{equation}
\label{eq:spike-init-rev}
\frac{1}{d}G_0G_0^\top
=
\zeta_0I_m+V_\star\Theta_\star V_\star^\top .
\end{equation}
\end{assumption}

\begin{lemma}[Concentration of a random subspace projection]
\label{lem:subspace-concentration-rev}
Let \(1\le r<m\). Let \(U\in\mathbb R^{m\times r}\) be an orthonormal basis
of a Haar-random \(r\)-dimensional subspace of \(\mathbb R^m\), and let
\(v\in\mathbb R^m\) be any fixed unit vector. Then
\[
Z:=\|U^\top v\|_2^2
\sim
\mathrm{Beta}\!\left(\frac{r}{2},\frac{m-r}{2}\right),
\qquad
\mathbb E Z=\frac{r}{m} .
\]
Moreover, for every \(0<\varepsilon\le 1/3\),
\[
\mathbb P\left(
\left|Z-\frac{r}{m}\right|>\varepsilon\frac{r}{m}
\right)
\le
2\exp\left(-\frac{r\varepsilon^2}{5.2}\right).
\]
\end{lemma}

\begin{lemma}[First-step GD stable rank under the rank-\(k\) spike model]
\label{lem:firststep-gd-rev}
Under \eqref{eq:init-firststep-rev} and Assumption~\ref{assump:spike-init-rev}, fix \(\delta\in(0,1)\) and \(u\ge0\). Define
\[
\varepsilon_\delta:=\sqrt{\frac{5.2\log(4k/\delta)}{r}},
\]
and assume \(\varepsilon_\delta\le 1/3\) and \(1-\sqrt{r/m}-u/\sqrt m>0\). Then, with probability at least \(1-\delta-2e^{-u^2/2}\) over the Gaussian initialization of \(A_0\),
\[
\srank(U_0^\top G_0)
\le
\frac{1+\varepsilon_\delta}{1-\varepsilon_\delta}\srank(G_0),
\]
and
\[
\srank(\Delta W_0^{\mathrm{gd}})
\le
\frac{1+\varepsilon_\delta}{1-\varepsilon_\delta}
\left(
\frac{1+\sqrt{r/m}+u/\sqrt m}{1-\sqrt{r/m}-u/\sqrt m}
\right)^4
\srank(G_0).
\]
\end{lemma}

\begin{proof}
At initialization,
\[
\Delta W_0^{\mathrm{gd}}\propto -A_0A_0^\top G_0,
\qquad
A_0A_0^\top=U_0\Sigma_0^2U_0^\top.
\]
Let
\[
H_0:=U_0^\top G_0,
\qquad
C_0:=\frac{1}{d}H_0H_0^\top
=
\frac{1}{d}U_0^\top G_0G_0^\top U_0.
\]
The stable rank of the GD update can be expressed through \(C_0\) and the singular values of \(A_0\):
\[
\|A_0A_0^\top G_0\|_F^2
=
d\tr(\Sigma_0^4C_0),
\qquad
\|A_0A_0^\top G_0\|_2^2
=
d\lambda_{\max}(\Sigma_0^2C_0\Sigma_0^2).
\]
Equivalently,
\[
\srank(\Delta W_0^{\mathrm{gd}})
=
\srank(\Sigma_0^2C_0^{1/2}).
\]
For any matrix \(X\),
\[
\sigma_r(\Sigma_0^2)\|X\|_F
\le
\|\Sigma_0^2X\|_F
\le
\sigma_1(\Sigma_0^2)\|X\|_F,
\]
and
\[
\sigma_r(\Sigma_0^2)\|X\|_2
\le
\|\Sigma_0^2X\|_2
\le
\sigma_1(\Sigma_0^2)\|X\|_2.
\]
Applying these inequalities with \(X=C_0^{1/2}\) gives
\[
\srank(\Delta W_0^{\mathrm{gd}})
\le
\left(\frac{\sigma_1(A_0)}{\sigma_r(A_0)}\right)^4
\srank(C_0^{1/2}).
\]
Since
\[
\srank(C_0^{1/2})
=
\frac{\tr(C_0)}{\lambda_{\max}(C_0)}
=
\srank(U_0^\top G_0),
\]
it remains to control \(\srank(U_0^\top G_0)\).

Under Assumption~\ref{assump:spike-init-rev},
\[
C_0=\zeta_0I_r+\Omega_0\Theta_\star\Omega_0^\top,
\qquad
\Omega_0:=U_0^\top V_\star.
\]
Thus
\[
\srank(U_0^\top G_0)
=
\frac{r\zeta_0+\tr(\Omega_0\Theta_\star\Omega_0^\top)}{\zeta_0+\lambda_{\max}(\Omega_0\Theta_\star\Omega_0^\top)}.
\]
For each spike direction \(v_j\), Lemma~\ref{lem:subspace-concentration-rev} gives
\[
\mathbb P\left(
\left|\|U_0^\top v_j\|_2^2-\frac{r}{m}\right|
>
\varepsilon\frac{r}{m}
\right)
\le
2\exp\left(-\frac{r\varepsilon^2}{5.2}\right).
\]
Taking \(\varepsilon=\varepsilon_\delta\) and applying a union bound over \(j=1,\ldots,k\), with probability at least \(1-\delta\),
\[
(1-\varepsilon_\delta)\frac{r}{m}
\le
\|U_0^\top v_j\|_2^2
\le
(1+\varepsilon_\delta)\frac{r}{m}
\qquad
\text{for all }j.
\]
On this event,
\[
\tr(\Omega_0\Theta_\star\Omega_0^\top)
=
\sum_{j=1}^k\theta_j\|U_0^\top v_j\|_2^2
\le
(1+\varepsilon_\delta)\frac{r}{m}\tr(\Theta_\star),
\]
and the top spike direction gives
\[
\lambda_{\max}(\Omega_0\Theta_\star\Omega_0^\top)
\ge
\theta_1\|U_0^\top v_1\|_2^2
\ge
(1-\varepsilon_\delta)\frac{r}{m}\theta_1.
\]
Substitution yields
\[
\srank(U_0^\top G_0)
\le
\frac{r\zeta_0+(1+\varepsilon_\delta)\frac{r}m\tr(\Theta_\star)}{\zeta_0+(1-\varepsilon_\delta)\frac{r}m\theta_1}.
\]
Writing \(\rho=r/m\), the right-hand side becomes
\[
\rho
\frac{m\zeta_0+(1+\varepsilon_\delta)\tr(\Theta_\star)}{\zeta_0+(1-\varepsilon_\delta)\rho\theta_1}.
\]
We compare this with
\[
\srank(G_0)
=
\frac{m\zeta_0+\tr(\Theta_\star)}{\zeta_0+\theta_1}.
\]
First,
\[
\frac{m\zeta_0+(1+\varepsilon_\delta)\tr(\Theta_\star)}{m\zeta_0+\tr(\Theta_\star)}
\le
1+\varepsilon_\delta.
\]
Second, since \(\rho\le1\),
\[
(1-\varepsilon_\delta)\rho(\zeta_0+\theta_1)
\le
\zeta_0+(1-\varepsilon_\delta)\rho\theta_1,
\]
and therefore
\[
\rho
\frac{\zeta_0+\theta_1}{\zeta_0+(1-\varepsilon_\delta)\rho\theta_1}
\le
\frac{1}{1-\varepsilon_\delta}.
\]
Combining the preceding displays gives
\[
\srank(U_0^\top G_0)
\le
\frac{1+\varepsilon_\delta}{1-\varepsilon_\delta}\srank(G_0).
\]

Finally, standard Gaussian singular-value bounds for \(A_0\) give, with probability at least \(1-2e^{-u^2/2}\),
\[
\sigma_1(A_0)\le1+\sqrt{r/m}+u/\sqrt m,
\qquad
\sigma_r(A_0)\ge1-\sqrt{r/m}-u/\sqrt m.
\]
Combining this event with the random-subspace concentration event proves the claimed GD bound.
\end{proof}

\end{document}
